\section{Theory freeze and entry criteria for code implementation}
\label{sec:theory-freeze}

This section marks the end of architecture design required before the first
implementation branch.

\subsection{Status}

The vNext theory is considered \textbf{implementation-ready} for the MVP in
Section~\ref{sec:mvp} when the following are defined:
\begin{enumerate}
\item continuous Scene and superelliptic conductor geometry;
\item admissibility and material-domain assumptions;
\item homogeneous current--potential--charge integral teacher;
\item canonical work-conjugate port definition;
\item singular/near/far quadrature policy;
\item low-frequency-stable mixed system and correctness solver;
\item finite-cross-section nested modal enrichment;
\item physical loss operator and exact power closure;
\item stable thermal transfer/reduced realization;
\item carrier-envelope electrothermal coupling;
\item physics-factored neural architecture;
\item FAST/CERTIFIED/REFERENCE semantics;
\item immutable data/split/active-learning policy;
\item analytic and independent benchmark suite;
\item explicit architecture falsification criteria.
\end{enumerate}
All of these are now specified by this document.

\subsection{Open questions that do not block coding}

The following remain research choices and are expected to be resolved by
prototype measurements:
\begin{enumerate}
\item optimal Jacobi/boundary-layer conductor mode ordering;
\item exact knot-refinement heuristic along the spiral;
\item FMM versus hierarchical-matrix acceleration;
\item best scalable Schur/preconditioner approximation;
\item graph architecture and latent widths;
\item best loss-channel joint-reduction algorithm;
\item uncertainty-calibration method;
\item broadband reciprocal port-Hamiltonian realization;
\item SIE/VIE formulations beyond the MVP.
\end{enumerate}
None changes the MVP public Scene, port, power, loss, thermal, or certification
semantics. Therefore they are not reasons to delay implementation.

\subsection{Implementation phases}

The first implementation branch proceeds in five phases.

\paragraph{I0: geometry and analytic benchmarks.}
Implement continuous superellipse/sweep geometry, Bishop frame, quadrature
patches, terminals, rigid-transform tests, and benchmark formulas. No neural
network.

\paragraph{I1: dense homogeneous conductor teacher.}
Implement the mixed KKT operator of Section~\ref{sec:stable-mixed}, dense
singular/near interaction blocks, port extraction, loss integration, and nested
basis/quadrature convergence. This phase is the first hard go/no-go point.

\paragraph{I2: thermal transfer and electrothermal loop.}
Implement loss channels, analytic/compact thermal transfer, stable reduced
thermal dynamics, carrier-envelope coupling, and time/steady benchmarks.

\paragraph{I3: physics-factored FAST surrogate.}
Generate immutable REFERENCE data, train self/pair/many-body structured models,
continuous loss decoder, and thermal correction. FAST must pass fixed validation.

\paragraph{I4: CERTIFIED correction and scale-up.}
Add matrix-free/hierarchical/FMM application, block FGMRES, FAST-to-physical
lifting, goal-oriented certification, active learning, and performance caches.

Package SIE/VIE and broadband EM begin only after I0--I4 demonstrate the MVP
architecture.

\subsection{Go/no-go after I1}

Continue to neural implementation only if the dense homogeneous conductor
teacher simultaneously demonstrates:
\begin{enumerate}
\item DC limit;
\item coaxial mutual inductance;
\item generic-pose Neumann thin-wire limit;
\item finite-cross-section skin/proximity convergence;
\item reciprocity;
\item rigid-motion invariance;
\item power closure;
\item basis and quadrature convergence without world-volume refinement.
\end{enumerate}

If these fail because of the chosen mixed formulation or modal space, revise
that isolated numerical choice before writing neural code. Do not change the
public architecture to rescue the teacher.

\subsection{Theory change control after code begins}

After the implementation branch starts:
\begin{enumerate}
\item corrections to algebra/signs/definitions are allowed with a theory
      version increment;
\item numerical basis/preconditioner choices may change without changing public
      schema if invariants are preserved;
\item changing Scene semantics, port work convention, continuous loss
      definition, FAST/CERTIFIED meaning, or removing hard physical invariants
      requires an explicit architecture revision;
\item implementation difficulties alone are not sufficient justification for
      restoring a global world-volume mesh.
\end{enumerate}

\subsection{Conclusion}

There are no remaining theoretical blockers to starting I0 and I1. Further
theoretical work should now be driven by concrete failures of the new
implementation or by post-MVP extensions, not by speculative expansion of the
architecture.
